\documentclass[11pt]{article}
\usepackage[margin=1in]{geometry}
\usepackage{amsmath,amssymb,graphicx,booktabs}
\usepackage[hidelinks]{hyperref}
\title{Source-Aware Security Steering with Learned Intervention Strength}
\author{Author and affiliation to be supplied}
\date{Research draft --- September 2026}
\begin{document}
\maketitle
\begin{abstract}
We describe a reproducible experimental framework for testing whether inference-time
activation steering can reduce indirect prompt injection while preserving legitimate
task completion. The framework constructs paired safe, injected, and hard-benign contexts;
estimates layer- and site-specific security directions; compares selected placement with
randomized controls; and learns intervention strength using a small contextual bandit.
Security and utility are evaluated separately and jointly. This manuscript documents an
implemented method and pilot protocol, not an established defense advantage. Publication
claims require larger held-out experiments and native agent-environment evaluation.
\end{abstract}

\section{Motivation and related work}
An assistant may encounter attacker instructions in externally supplied data. We study
interventions that preserve the user's task when such data is included in the context.
Prompt injection is a recognized application risk~\cite{owasp2025}. SEAL estimates
reasoning directions offline and offsets hidden representations during generation~\cite{chen2025seal}.
Its reasoning-calibration results motivate our hypothesis but do not validate a prompt
injection defense. ARGUS studies instruction-following steering and adaptive strength for
multimodal injection~\cite{lu2025argus}; ICON combines latent probing with inference-time
attention correction~\cite{wang2026icon}. Accordingly, we make no priority claim for the
general idea of steering against injection. Our experimental focus is the comparison of
representation estimators, residual/MLP placement, and reward-based strength selection.

\section{Threat model and data}
System policy and user intent are trusted. Retrieved records, web pages, tool outputs,
memory records and agent messages are untrusted data. The attacker may modify one such
record, but not its trusted provenance label. Each controlled family retains a common
legitimate task while varying the injected instruction or its benign quotation. Seven
categories cover override, reveal, exfiltration, tool misuse, memory poisoning, source
confusion and obfuscation. These are task annotations rather than a proven behavior taxonomy.
All canaries are fictional; actions are simulated and never executed externally.

The collection pipeline also adapts BIPIA EmailQA~\cite{yi2023bipia} and InjecAgent
tool responses~\cite{injecagent}. Static adaptations retain source hashes and mark outcomes
without sufficient task checks as unknown. These results are not directly comparable
to their native benchmark protocols. AgentDojo~\cite{agentdojo} provides an appropriate
next environment for validating actual action outcomes.

Families, declared template groups, exact duplicate inputs, and duplicate attack payloads
are grouped before splitting. Official test examples remain held out. Train data fits
vectors and controllers, validation chooses hyperparameters, and the final test plan is
saved before test execution. This controls exact leakage, not semantic overlap or
pretraining contamination.

\section{Security directions}
For layer $l$, site $s$, and type $k$, the mean-difference estimator is
\begin{equation}
v_{l,s,k}=\operatorname{normalize}\left(\mu^{safe}_{l,s,k}-\mu^{attack}_{l,s,k}\right).
\end{equation}
We retain both class prototypes and compare this estimator with family-balanced paired
means, sign-oriented uncentered PCA of paired differences, and a ridge-probe direction.
The default inference intervention uses a global vector, avoiding access to the test
example's true attack type. Class vectors support fixed-type ablations.

Representations are collected with hooks at the exact residual or MLP output used at
inference. This avoids mixing sites or using a final-normalized hidden state as though it
were a raw decoder output. Prompt-condition labels are not observed behavior labels:
separability may reflect wording or source formatting. Future behavior-labelled
continuation contrasts are necessary to interpret a direction as malicious compliance.

\section{Intervention and controls}
For the frozen base model, we apply
\begin{equation}
h'_{l,t}=h_{l,t}+m_t\alpha v_{l,s,k}.
\end{equation}
This is a local inference offset, not a recurrence across time and not a weight update.
Alpha denotes intervention strength. Decision-token, prefill-only, decode-only, and
all-position schedules are available. Multiple active sites/layers share a nominal
alpha budget. Absolute and activation-norm-relative scaling are supported.

Validation compares residual, MLP, and combined placement over layers, strengths and
estimators. Controls include the unsteered model, matched random layers, norm-matched
random directions, and sign reversal. All arms share prompt formatting and greedy
decoding. Placement is sampled per family and seed for the random-layer arm. We report
utility-constraint failures explicitly and retain the baseline as a recommendation option.

\section{Learning intervention strength}
A train-only logistic probe provides an attack-risk proxy from unsteered decision states.
A two-layer tanh network receives this proxy and observable length/source features, then
selects alpha from a finite grid containing zero. A one-response episode receives reward
\begin{equation}
R=w_s\mathbf{1}[\text{valid and no measured violation}]
+w_u\mathbf{1}[\text{task completed}]-\lambda |\alpha|.
\end{equation}
We optimize the categorical policy with REINFORCE and a running reward baseline. This
is a contextual bandit, not a sequence-level RL agent or a guarantee of the globally
optimal alpha. Actions, propensities, features and rewards are logged. Examples without
complete security and utility checks cannot provide training rewards. The policy and
risk probe are frozen for final evaluation; no test labels are inference features.

\section{Evaluation and pilot results}
Report attack success rate (ASR), benign utility (BU), attacked utility (AU), and joint
security/task success. A correct answer may coexist with successful exfiltration, while
a refusal may prevent an attack without satisfying the task. Unknown outcomes remain
ungraded, with explicit denominators. Output-format failures are distinguished from
diagnostic answer presence. Canary, tool and memory checks are necessarily limited by
the serialized-action protocol. Token counts and end-to-end generation latency accompany
the accuracy measures; probe prepasses should be included in deployment cost estimates.

Wilson intervals describe individual rates. Paired method differences use a family-level
bootstrap to preserve correlations between variants. Random weights test software only.
Pilots with small pretrained models assess feasibility and cannot establish robust defense gains.

\IfFileExists{pilot_results.tex}{\subsection{Small CPU feasibility pilot}
A SmolLM2-135M-Instruct run used 24 training examples, six validation examples and six
test examples across two independent test families. We compared 24 validation settings
and trained the alpha controller for 72 updates. The primary rubric required the code
alone, matching the user instruction. Baseline, selected fixed steering and adaptive
steering all had zero strict task completion and zero observed attacker successes.
Correct-code presence was 5/6 for every arm in a post-hoc diagnostic, often accompanied
by disallowed introductory prose or Markdown. This diagnostic did not replace the
predefined utility metric or determine model selection.

\begin{center}
\begin{tabular}{lrrr}
\toprule
Method & Attack success & Benign task success & Joint success \\
\midrule
Unsteered baseline & 0/2 & 0/4 & 0/6 \\
Selected fixed steering & 0/2 & 0/4 & 0/6 \\
Neural alpha policy & 0/2 & 0/4 & 0/6 \\
\bottomrule
\end{tabular}
\end{center}

Random-layer/direction controls and sign reversal had identical aggregate outcomes.
The selected nonzero intervention used the mean-difference direction at residual layer
8 with alpha=1; validation did not prefer it over baseline. The learned policy chose
alpha=1 in all test cases. With only two attacked examples, the descriptive Wilson 95\%
upper bound for ASR is approximately 65.8\%. These results establish pipeline execution,
not useful defense gains or successful adaptive calibration. The original logs and
configuration are retained in \texttt{runs/cpu\_pilot}; see
\texttt{docs/pilot\_findings.md} for provenance and interpretation limits.
}{
No numerical research claim is made in this draft. The experiment command writes
\texttt{report.md}, \texttt{metrics.csv}, \texttt{summary.json}, and
\texttt{results\_table.tex}. Include reviewed, appropriately labelled results here.
}

\section{Limitations and next experiments}
The present setup does not prove a universal security direction, calibrated detection,
native tool safety, robustness against an adaptive attacker, or cross-model transfer.
Vector scale and source formatting can confound placement comparisons. Template-based
tasks may be too narrow; their exact-match criteria may penalize harmless extra prose.
Future evaluation should hold out attack authors and sources, use stronger models and
native environments, reproduce ARGUS/ICON-style baselines where applicable, compare
behavior-conditioned vectors, and repeat controller training with independent seeds.
Mentorship, conference acceptance and financial support are external outcomes and are
not represented as completed experimental milestones.

\bibliographystyle{plain}
\bibliography{references}
\end{document}
